% !TEX root = ../main.tex

\section{Kwantowy silnik Stirlinga}
W kwantowym silniku Stirlinga pracującym równowagowo ciepło dostarczone podczas przemian izotermicznych możemy znaleźć wykorzystując wzór \eqref{eq:def_entropia}, natomiast w przemianach izochorycznych jest to różnica energii wewnętrznej zgodnie z Pierwszą Zasadą Termodynamiki (\(W=0\)). Ciepła te są równe:
\begin{align*}
	Q_{1-2} &= S\pab{T_2,\omega_2} - S\pab{T_1,\omega_1}, \\
	Q_{2-3} &= U\pab{T_3,\omega_2} - U\pab{T_2,\omega_2}, \\
	Q_{3-4} &= S\pab{T_4,\omega_1} - S\pab{T_3,\omega_2}, \\
	Q_{4-1} &= U\pab{T_1,\omega_1} - U\pab{T_4,\omega_2},
\end{align*}
gdzie \(\omega\) jest parametrem kontrolowanym z zewnątrz i ma sens odwrotnej objętości \cite{myers}. Dokładna definicja \(\omega\) zależna jest od rozpatrywanego czynnika roboczego.

\subsection{Oscylator harmoniczny}
Czynnikiem roboczym jest oscylator harmoniczny, a parametr \(\omega\) jest częstością drgań własnych oscylatora. W dodatku \ref{sec:osc_harm} wyprowadzono wielkości termodynamiczne U \ref{eq:osc_harm_U} i S \ref{eq:osc_harm_S}. Wykorzystując je znajdujemy odpowiednie ciepła w przemianach izotermicznych:
\begin{gather}
\begin{split}
	Q_{1-2} &= T_c S\pab{T_c,\omega_2} - T_c S\pab{T_c,\omega_1} = \\ &=
	kT_c \ln\bab{2\sinh\pab{\frac{\hbar\omega_1}{2k_B T_c}}} -
	kT_c \ln\bab{2\sinh\pab{\frac{\hbar\omega_2}{2k_B T_c}}} + \\ &+
	\frac{\hbar\omega_2}{2} \ctgh\pab{\frac{\hbar\omega_2}{2k_B T_c}} -
	\frac{\hbar\omega_1}{2} \ctgh\pab{\frac{\hbar\omega_1}{2k_B T_c}},
\end{split} \\
\begin{split}
	Q_{3-4} &= T_h S\pab{T_h,\omega_1} - T_h S\pab{T_h,\omega_2} = \\ &=
	kT_h \ln\bab{2\sinh\pab{\frac{\hbar\omega_2}{2k_B T_h}}} -
	kT_h \ln\bab{2\sinh\pab{\frac{\hbar\omega_1}{2k_B T_h}}} + \\ &+
	\frac{\hbar\omega_1}{2} \ctgh\pab{\frac{\hbar\omega_1}{2k_B T_h}} -
	\frac{\hbar\omega_2}{2} \ctgh\pab{\frac{\hbar\omega_2}{2k_B T_h}}
\end{split}
\end{gather}
oraz w przemianach izochorycznych:
\begin{gather}
\begin{split}
	Q_{2-3} &= U\pab{T_h,\omega_2} - U\pab{T_c,\omega_2} = \\ &=
	\frac{\hbar\omega_2}{2}\ctgh\pab{\frac{\hbar\omega_2}{2k_B T_h}} -
	\frac{\hbar\omega_2}{2}\ctgh\pab{\frac{\hbar\omega_2}{2k_B T_c}},
\end{split} \\
\begin{split}
	Q_{4-1} &= U\pab{T_c,\omega_1} - U\pab{T_h,\omega_1} = \\ &=
	\frac{\hbar\omega_1}{2}\ctgh\pab{\frac{\hbar\omega_1}{2k_B T_c}} -
	\frac{\hbar\omega_1}{2}\ctgh\pab{\frac{\hbar\omega_1}{2k_B T_h}}.
\end{split}
\end{gather}

W przeciwieństwie do klasycznego silnika Stirlinga ciepła te nie są sobie równe co do wielkości i nie znoszą się nawzajem. Człony te jednak skracają się z częścią ciepeł w przemianach izotermicznych. Zgodnie ze wzorem \eqref{eq:praca_w_silniku} znajdujemy pracę:
\begin{equation}
\begin{split}
	- W &= Q_{1-2} + Q_{2-3} + Q_{3-4} + Q_{4-1} = \\ =
	kT_h \ln\bab{\frac{
	2\sinh\pab{\frac{\hbar\omega_2}{2k_B T_h}}
	}{
	2\sinh\pab{\frac{\hbar\omega_1}{2k_B T_h}}
	}} +
	kT_c \ln\bab{\frac{
	2\sinh\pab{\frac{\hbar\omega_1}{2k_B T_c}}
	}{
	2\sinh\pab{\frac{\hbar\omega_2}{2k_B T_c}}
	}}.
\end{split}
\end{equation}

Oznaczamy ciepło wpływające z regeneratora do czynnika roboczego poprzez \cite{myers}:
\begin{equation}
	\Delta Q = Q_{2-3} + Q_{4-1}.
\end{equation}
W silniku klasycznym \(\Delta Q = 0\). W kwantowym jest to prawdziwe jedynie dla pewnych wartości \(\omega_1\), \(\omega_2\), \(T_c\) oraz \(T_h\). Jeśli \(\Delta Q < 0\) to więcej ciepła wpływa do regeneratora niż z niego wypływa i nadmiarowe ciepło musi być odprowadzone do zbiornika zimnego. Jeśli \(\Delta Q > 0\) to zachodzi sytuacja odwrotna i dodatkowe ciepło musi być przez regenerator pobrane ze zbiornika ciepłego.

\begin{figure}[h]
	\centering
	\begin{subfigure}{0.55\textwidth}
		\centering
		\begin{tikzpicture}
	\begin{axis}[xlabel=$\omega_1$,
		ylabel=$\frac{\omega_2}{\omega_1}$,
		zlabel=$\eta$,
		grid=major,
		mesh/ordering=y varies,
		mesh/cols=51,
		mesh/rows=51,
		view={-45}{45}]
		\addplot3[surf] file {Program/Stirling_harm_osc_data.dat};
	\end{axis}
\end{tikzpicture}
	\end{subfigure}
	\begin{subfigure}{0.55\textwidth}
		\centering
		\begin{tikzpicture}
	\begin{axis}[xlabel=$\omega_1$,
		ylabel=$\kappa$,
		zlabel=$\eta$,
		grid=major,
		colorbar,
		mesh/ordering=y varies,
		mesh/cols=51,
		mesh/rows=51,
		view={0}{90}]
		\addplot3[surf] file {Program/Stirling_harm_osc_data.dat};
	\end{axis}
\end{tikzpicture}
	\end{subfigure}
	\caption{Wykres sprawności kwantowego silnika Stirlinga z oscylatorem harmonicznym jako czynnikiem roboczym}
	\label{fig:wykres_kw_stirling_osc_harm}
\end{figure}

\subsection{Układ o spinie \(\inv{2}\)}
Przemiany izotermiczne:
\begin{gather}
\begin{split}
	Q_{1-2} &= T_c S\pab{T_c,\omega_2} - T_c S\pab{T_c,\omega_1} = \\ &=
	kT_c \ln\bab{2\cosh\pab{\frac{\hbar\omega_2}{2k_B T_c}}} -
	kT_c \ln\bab{2\cosh\pab{\frac{\hbar\omega_1}{2k_B T_c}}} + \\ &+
	\frac{\hbar\omega_1}{2} \tgh\pab{\frac{\hbar\omega_1}{2k_B T_c}} -
	\frac{\hbar\omega_2}{2} \tgh\pab{\frac{\hbar\omega_2}{2k_B T_c}},
\end{split} \\
\begin{split}
	Q_{3-4} &= T_h S\pab{T_h,\omega_1} - T_h S\pab{T_h,\omega_2} = \\ &=
	kT_h \ln\bab{2\cosh\pab{\frac{\hbar\omega_1}{2k_B T_h}}} -
	kT_h \ln\bab{2\cosh\pab{\frac{\hbar\omega_2}{2k_B T_h}}} + \\ &+
	\frac{\hbar\omega_2}{2} \tgh\pab{\frac{\hbar\omega_2}{2k_B T_h}} -
	\frac{\hbar\omega_1}{2} \tgh\pab{\frac{\hbar\omega_1}{2k_B T_h}}.
\end{split}
\end{gather}

Przemiany izochoryczne:
\begin{gather}
\begin{split}
	Q_{2-3} &= U\pab{T_h,\omega_2} - U\pab{T_c,\omega_2} = \\ &=
	\frac{\hbar\omega_2}{2}\tgh\pab{\frac{\hbar\omega_2}{2k_B T_c}} -
	\frac{\hbar\omega_2}{2}\tgh\pab{\frac{\hbar\omega_2}{2k_B T_h}}
\end{split} \\
\begin{split}
	Q_{4-1} &= U\pab{T_c,\omega_1} - U\pab{T_h,\omega_1} = \\ &=
	\frac{\hbar\omega_1}{2}\tgh\pab{\frac{\hbar\omega_1}{2k_B T_h}} -
	\frac{\hbar\omega_1}{2}\tgh\pab{\frac{\hbar\omega_1}{2k_B T_c}}
\end{split}
\end{gather}

Praca cyklu jest równa:
\begin{equation}
\begin{split}
	-W &= Q_{1-2} + Q_{2-3} + Q_{3-4} + Q_{4-1} = \\ &=
	kT_h \ln\bab{\frac{
	\cosh\pab{\frac{\hbar\omega_1}{2k_B T_h}}
	}{
	\cosh\pab{\frac{\hbar\omega_2}{2k_B T_h}}
	}} +
	kT_c \ln\bab{\frac{
	\cosh\pab{\frac{\hbar\omega_2}{2k_B T_c}}
	}{
	\cosh\pab{\frac{\hbar\omega_1}{2k_B T_c}}
	}}
\end{split}
\end{equation}

\begin{figure}[h]
	\centering
	\begin{subfigure}{0.55\textwidth}
		\centering
		\begin{tikzpicture}
	\begin{axis}[xlabel=$\omega_1$,
		ylabel=$\kappa$,
		zlabel=$\eta$,
		grid=major,
		mesh/ordering=y varies,
		mesh/cols=51,
		mesh/rows=51,
		view={-45}{45}]
		\addplot3[surf] file {Program/Stirling_spin_pol_data.dat};
	\end{axis}
\end{tikzpicture}
	\end{subfigure}
	\begin{subfigure}{0.55\textwidth}
		\centering
		\begin{tikzpicture}
	\begin{axis}[xlabel=$\omega_1$,
		ylabel=$\frac{\omega_2}{\omega_1}$,
		zlabel=$\eta$,
		grid=major,
		colorbar,
		mesh/ordering=y varies,
		mesh/cols=51,
		mesh/rows=51,
		view={0}{90}]
		\addplot3[surf] file {Program/Stirling_spin_pol_data.dat};
	\end{axis}
\end{tikzpicture}
	\end{subfigure}
	\caption{Wykres sprawności kwantowego silnika Stirlinga z układem o spinie \(\inv{2}\) jako czynniku roboczym}
	\label{fig:wykres_kw_stirling_spin_pol}
\end{figure}